\documentclass[../../main.tex]{subfiles}
\graphicspath{{../../image/}{../../figure/}}
\begin{document}

\section{熵激波容许性准则}\label{section:ch08-熵激波容许性准则}

使用熵不等式来剔除一般双曲守恒律系统的虚假弱解这一思想，是在\cref{section:ch04-熵可容性条件}中引入的，并在第 IV、V、VI 章中反复使用。我们观察到，在 $BV$ 弱解的语境下，熵条件蜕化为在激波集每一点处检验的不等式组\eqref{align:ch04-拟凸}。对一维系统\eqref{align:ch08-守恒律一维}，\eqref{align:ch04-拟凸} 取形式
\begin{align}
-s\big[\eta(U_+)-\eta(U_-)\big]+q(U_+)-q(U_-)\leqslant0,
\label{align:ch08-熵条件}
\end{align}
其中 $(\eta,q)$ 是满足\eqref{align:ch07-熵流对定义} 的熵--熵流对，$\mathrm{D}q=\mathrm{D}\eta\,\mathrm{D}F$。\eqref{align:ch08-熵条件} 左端的量今后称为跨过激波的\emph{熵产生}。

熵条件蜕化为对激波的逐点检验这一事实，在形成下述主流观点方面起了主导作用：容许性只需在激波层面检验，即一般 $BV$ 弱解是容许的，当且仅当其每一条激波都是容许的。

在建立熵容许性条件\eqref{align:ch08-熵条件} 时，首要任务是指定合适的熵--熵流对 $(\eta,q)$。每当\eqref{align:ch08-守恒律一维} 与物理相联系时，总应指定物理上合适的熵。特别地，对系统\eqref{align:ch07-热弹性}、\eqref{align:ch07-p系统}、\eqref{align:ch07-理想气体}，应分别指定熵对\eqref{align:ch07-热弹性熵}、\eqref{align:ch07-p系统熵}、\eqref{align:ch07-理想气体熵}。\footnote{在将\eqref{align:ch08-熵条件} 应用于配有熵--熵流对\eqref{align:ch07-热弹性熵} 的系统 (7.1.8) 时，不应混淆 (7.4.9) 中的 $s$（即物理熵）与\eqref{align:ch08-熵条件} 中的 $s$（激波速度）。由于 $q=0$，此处\eqref{align:ch08-熵条件} 表述为“激波通过后物理熵必须增加”。提请读者注意，这一命题有时被误解为一条普适物理原理，并在其不再适用时仍被套用。}

在缺乏物理指导，或物理所提供的熵--熵流对不足以剔除所有虚假激波时，必须（在可行时）指定由其他容许性准则（如黏性）驱动的附加熵--熵流对。在此方面，应记住：如前几章所示，熵函数的凸性是一个理想特征。

我们从单个守恒律\eqref{align:ch07-标量守恒律} 开始考察。激波速度 $s$ 由\eqref{align:ch08-扰动方程} 给出。依照第 VI 章的讨论，容许激波必须对所有凸函数 $\eta$ 满足\eqref{align:ch08-熵条件}。然而，如\cref{section:ch06-可容许弱解稳定性}所解释，\eqref{align:ch08-熵条件} 只需对熵--熵流对族\eqref{align:ch06-6-2-5} 检验，即
\begin{align}
\eta(u;\bar u)=(u-\bar u)^+,\qquad q(u;\bar u)=\operatorname{sgn}(u-\bar u)^+\big[f(u)-f(\bar u)\big].
\label{align:ch08-标量熵族}
\end{align}
立即可见：对族\eqref{align:ch08-标量熵族} 中每个 $(\eta,q)$，\eqref{align:ch08-熵条件} 成立，当且仅当对介于 $u_-$ 与 $u_+$ 之间的所有 $u_0$，\eqref{align:ch08-oleinik} 成立。我们由此重新推回了在\hyperref[section:ch08-liu容许性准则]{8.4 节}中遇到的 Oleinik $E$-条件。这意味着，对单个守恒律，对所有凸熵施加的熵容许性条件等价于 Liu $E$-条件。

一般而言，不可能由单个熵--熵流对的熵条件\eqref{align:ch08-熵条件} 恢复 Oleinik $E$-条件。例如取
\begin{align}
\eta(u)=\tfrac12 u^2,\qquad q(u)=\int_0^u \omega f'(\omega)\,\mathrm{d}\omega.
\label{align:ch08-二次熵}
\end{align}
由\eqref{align:ch08-扰动方程} 并经简短计算，\eqref{align:ch08-熵条件} 取形式
\begin{align}
\tfrac12\big[f(u_+)+f(u_-)\big](u_+-u_-)-\int_{u_-}^{u_+}f(\omega)\,\mathrm{d}\omega\leqslant0.
\label{align:ch08-二次熵条件}
\end{align}
注意，此处跨过激波的熵产生由下述有向面积度量：以 $f$ 的图像端点 $(u_-,f(u_-))$、$(u_+,f(u_+))$ 为弧的图形与连接这两点的弦所围区域。显然，Oleinik $E$-条件\eqref{align:ch08-oleinik} 蕴涵\eqref{align:ch08-二次熵条件}，但逆命题一般不成立。进而，\eqref{align:ch08-二次熵条件} 一般既不蕴涵 Lax $E$-条件\eqref{align:ch08-标量lax}，逆命题也不成立。然而，当 $f$ 凸或凹时，\eqref{align:ch08-二次熵条件}、\eqref{align:ch08-oleinik}、\eqref{align:ch08-标量lax} 彼此等价。

接下来转向等熵弹性系统\eqref{align:ch07-p系统}。我们使用由\eqref{align:ch07-p系统熵} 给出的熵--熵流对 $(\eta,q)$。一个有趣而相当冗长的计算（涉及 Rankine--Hugoniot 条件\eqref{align:ch08-弹性rh}）表明，此处\eqref{align:ch08-熵条件} 蜕化为
\begin{align}
s\left\{\tfrac12\big[\sigma(u_+)+\sigma(u_-)\big](u_+-u_-)-\int_{u_-}^{u_+}\sigma(\omega)\,\mathrm{d}\omega\right\}\leqslant0.
\label{align:ch08-弹性熵条件}
\end{align}
\eqref{align:ch08-弹性熵条件} 左端花括号中的量，度量由 $\sigma$ 的图像端点 $(u_-,\sigma(u_-))$、$(u_+,\sigma(u_+))$ 弧与连接这两点的弦所围集合的有向面积。因此，Wendroff $E$-条件\eqref{align:ch08-wendroff} 蕴涵\eqref{align:ch08-弹性熵条件}，但逆命题一般不成立。条件\eqref{align:ch08-弹性熵条件} 未必蕴涵 Lax $E$-条件\eqref{align:ch08-弹性lax}，逆命题也不必然成立。然而，当 $\sigma$ 凸或凹时，\eqref{align:ch08-弹性熵条件}、\eqref{align:ch08-wendroff}、\eqref{align:ch08-弹性lax} 彼此等价。当然，系统\eqref{align:ch07-p系统} 配有丰富的熵集，故可使用附加熵--熵流对从熵条件恢复 Wendroff $E$-条件，但此处不作尝试。

我们现在对一般系统\eqref{align:ch08-守恒律一维}、在 $U_-$ 与 $U_+$ 由一条激波曲线相连的假设下，考虑熵激波容许性条件\eqref{align:ch08-熵条件}。特别地，这将涵盖中等强度激波的情形。于是设 $\lambda_i(U_-)$ 是简单特征速度，考虑过 $U_-$ 的 $i$-激波曲线 $W_i(\tau;U_-)$，并令 $U_+=W_i(\tau_+;U_-)$，$s=s_i(\tau_+;U_-)$。沿 $i$-激波曲线的熵产生由下式给出：
\begin{align}
E(\tau)=-s_i(\tau)\big[\eta(W_i(\tau))-\eta(U_-)\big]+q(W_i(\tau))-q(U_-).
\label{align:ch08-熵产生}
\end{align}
对\eqref{align:ch08-熵产生} 求导并使用\eqref{align:ch07-熵流对定义}，得
\begin{align}
\dot E=-\dot s_i\big[\eta(W_i)-\eta(U_-)\big]-s_i\,\mathrm{D}\eta(W_i)\dot W_i+\mathrm{D}\eta(W_i)\mathrm{D}F(W_i)\dot W_i.
\label{align:ch08-熵产生导}
\end{align}
将\eqref{align:ch08-熵产生导} 与\eqref{align:ch08-drh曲线}（取 $\bar U=U_-$）联立，推得
\begin{align}
\dot E=-\dot s_i\big\{\eta(W_i)-\eta(U_-)-\mathrm{D}\eta(W_i)\big[W_i-U_-\big]\big\}.
\label{align:ch08-熵产生导简}
\end{align}
注意\eqref{align:ch08-熵产生导简} 右端是激波强度的二阶小量。因此，跨过激波的熵产生 $E(\tau_+)$（即 $\dot E(\tau)$ 从 $0$ 到 $\tau_+$ 的积分）是 $\tau_+$ 的三阶小量。我们由此建立了下述

\begin{theorem}[弱激波熵产生的阶]\label{theorem:ch08-弱激波熵产生}
跨过弱激波的熵产生是激波强度的三阶小量。
\end{theorem}

当 $U_-$ 是 $i$-特征族的线性退化点时，$\dot s_i(0;U_-)$ 为零，于是跨过激波的熵产生将是激波强度的（至多）四阶小量。特别地，当 $i$-特征族线性退化时，由\cref{theorem:ch08-线性退化激波曲线}，$\dot s_i$ 恒为零，于是

\begin{theorem}[接触间断的熵产生为零]\label{theorem:ch08-接触熵零}
当 $i$-特征族线性退化时，跨过任一 $i$-激波（$i$-接触间断）的熵产生为零。
\end{theorem}

现在转向激波容许性问题。我们观察到：当 $\eta$ 为凸函数时，\eqref{align:ch08-熵产生导简} 右端花括号中的表达式非正。于是 $\dot E$ 与 $\dot s_i$ 同号。因此，若 $\tau_+<0$ 且 $\dot s_i\geqslant0$ 在 $(\tau_+,0)$ 上，或 $\tau_+>0$ 且 $\dot s_i\leqslant0$ 在 $(0,\tau_+)$ 上，则熵容许性条件 $E(\tau_+)\leqslant0$ 成立；而若 $\tau_+<0$ 且 $\dot s_i<0$ 在 $(\tau_+,0)$ 上，或 $\tau_+>0$ 且 $\dot s_i>0$ 在 $(0,\tau_+)$ 上，则被破坏。回顾\hyperref[section:ch08-lax容许性准则]{8.3 节}的讨论，我们得出：在 $\dot s_i(\tau)$ 不变号的、$0$ 任一侧的 $\tau$ 范围内，熵容许性条件与 Lax $E$-条件等价。特别地，当特征族真非线性且激波为弱激波时即如此：

\begin{theorem}[真非线性弱激波的熵与 Lax 条件等价]\label{theorem:ch08-熵lax等价}
当 $i$-特征族真非线性且 $\lambda_i$ 是简单特征速度时，弱 $i$-激波的熵容许性条件与 Lax $E$-条件等价。
\end{theorem}

为摆脱真非线性与弱激波的范围，考虑条件
\begin{align}
\tau\,\dot W_i^{\top}(\tau;U_-)\,\mathrm{D}^2\eta\big(W_i(\tau;U_-)\big)\big[W_i(\tau;U_-)-U_-\big]\geqslant0.
\label{align:ch08-熵二阶条件}
\end{align}
回顾\eqref{align:ch07-熵可积正交}、\eqref{align:ch07-熵凸} 与\cref{theorem:ch08-激波曲线}，我们得出：当熵 $\eta$ 凸时，\eqref{align:ch08-熵二阶条件} 对弱 $i$-激波恒成立；当 $i$-激波曲线延拓到中等强度激波范围时，它也将被满足；甚至对强激波也可能成立，只要 $\dot W_i$ 与 $W_i-U_-$ 几乎保持指向 $R_i$ 的方向。

\begin{theorem}[Liu 条件蕴涵熵条件]\label{theorem:ch08-liu蕴涵熵}
设过 $U_-$ 的 $i$-激波曲线 $W_i(\tau;U_-)$ 及相应激波速度函数 $s_i(\tau;U_-)$ 定义在包含 $0$ 的区间 $(\alpha,\beta)$ 上，并对 $\tau\in(\alpha,\beta)$ 满足\eqref{align:ch08-熵二阶条件}，其中 $\eta$ 是系统的一个凸熵。则任一在左侧连接 $U_-$、在右侧连接 $U_+=W_i(\tau_+;U_-)$、速度为 $s$ 且满足 Liu $E$-条件\eqref{align:ch08-liu} 的 $i$-激波，也满足熵容许性条件\eqref{align:ch08-熵条件}。
\end{theorem}

\begin{proof}
令
\begin{align}
Q(\tau)=\eta\big(W_i(\tau;U_-)\big)-\eta(U_-)-\mathrm{D}\eta\big(W_i(\tau;U_-)\big)\big[W_i(\tau;U_-)-U_-\big].
\label{align:ch08-Q}
\end{align}
由\eqref{align:ch08-熵二阶条件}，
\begin{align}
\tau\,\dot Q(\tau)\leqslant0.
\label{align:ch08-Q符号}
\end{align}
将\eqref{align:ch08-熵产生导简} 从 $0$ 到 $\tau_+$ 积分，分部积分，并使用\eqref{align:ch08-Q}、\eqref{align:ch08-Q符号} 与\eqref{align:ch08-liu}，得
\begin{align}
E(\tau_+)
&=-\int_0^{\tau_+}\dot s_i(\tau;U_-)Q(\tau)\,\mathrm{d}\tau
=-sQ(\tau_+)+\int_0^{\tau_+}s_i(\tau;U_-)\dot Q(\tau)\,\mathrm{d}\tau
\leqslant -sQ(\tau_+)+s\int_0^{\tau_+}\dot Q(\tau)\,\mathrm{d}\tau=0,
\label{align:ch08-熵产生积分}
\end{align}
这表明激波满足\eqref{align:ch08-熵条件}。证明完毕。
\end{proof}

在强激波领域，上述各种容许性准则的层级可能被破坏。我们已看到，对由真非线性条件 $\sigma''(u)<0$ 刻画的等熵热弹性系统\eqref{align:ch07-p系统}，任意强度的压缩激波都满足由热力学第二定律诱导的熵容许性准则\eqref{align:ch08-弹性熵条件}，以及 Liu $E$-条件\eqref{align:ch08-wendroff}。我们进而检验这是否也适用于直线绝热气流系统\eqref{align:ch07-热弹性}，即在第一与第三特征族真非线性条件 $p_{uu}>0$ 之下。

首先设状态方程在其定义域上满足 $2\theta-u\theta_u>0$，理想气体\eqref{align:ch02-2-5-20} 即是如此。于是控制激波曲线的微分方程\eqref{align:ch08-sde} 的初值问题是适定的。考虑状态 $(u_-,v_-,s_-)$ 与 $(u_+,v_+,s_+)$ 由一个以速度 $r_+>0$ 传播、压缩介质（即 $u_+>u_-$）的 3-激波相连。于是 $s_+=s(u_+)$，其中 $s(u)$ 表示\eqref{align:ch08-sde} 以 $\bar u=u_-$、初始条件 $s(u_-)=s_-$ 的解。由\eqref{align:ch08-sde} 与\eqref{align:ch08-sde2} 推知 $\mathrm{d}s/\mathrm{d}u$ 与 $\mathrm{d}^2s/\mathrm{d}u^2$ 在 $u_-$ 处均为零。进而，对\eqref{align:ch08-sde2} 关于 $u$ 求导，得 $\mathrm{d}^3s/\mathrm{d}u^3<0$ 在 $u_-$ 处。因此，至少对 $u-u_-$ 为正且小，$\mathrm{d}s/\mathrm{d}u<0$。上述与\cref{theorem:ch08-弱激波熵产生}、\cref{theorem:ch08-熵lax等价}、\cref{theorem:ch08-liu蕴涵熵} 一致。我们现在注意到：由\eqref{align:ch08-sde2}，对 $u>u_-$，$s(u)$ 的任一临界点必为严格极大值。因而对所有 $u\in(u_-,u_+)$，$\mathrm{d}s/\mathrm{d}u<0$。特别地，这蕴涵由熵--熵流对 $(-s,0)$ 诱导、表达热力学第二定律的熵激波容许性条件 $s_+<s_-$ 被该激波满足。进而由\eqref{align:ch08-rdr} 知激波速度 $r(u)$ 在 $(u_-,u_+)$ 上递减，故
\begin{align}
\sqrt{-p_u(u_-,s_-)}=r(u_-)>r(u)>r(u_+)=r_+,\qquad u_-<u<u_+.
\label{align:ch08-速度不等式}
\end{align}
类似地，$s_-=s(u_-)$，其中 $s(u)$ 现在表示\eqref{align:ch08-sde} 以 $\bar u=u_+$、初始条件 $s(u_+)=s_+$ 的解。类似论证给出 $\sqrt{-p_u(u_+,s_+)}<r_+$。我们得出：在 $p_{uu}>0$ 与 $2\theta-u\theta_u>0$ 的假设下，压缩介质的任意强度 3-激波均满足 Lax $E$-条件、Liu $E$-条件与熵激波容许性条件。压缩 1-激波同样如此。

与之相反，可以构造满足热力学关系\eqref{align:ch07-热力学} 与双曲性条件\eqref{align:ch07-热弹性双曲}、但允许 $2\theta-u\theta_u$ 变号的状态方程，对它们系统\eqref{align:ch07-热弹性} 支撑满足 Liu $E$-条件但违反热力学第二定律的压缩激波。此类条件可在气体发生相变时出现。

\end{document}
